% !PRINTED 346
\weberpage{346}
%<<CONT
also wenn $\xi_1 + \xi_2 + \xi_3$ und $\eta_1 + \eta_2 + \eta_3$ mit Null congruent
ist, so gilt dasselbe von den linken Seiten von (7).

Nun gehen auch umgekehrt durch jede lineare Substitution der Form (6) immer drei
Grössen ($\xi, \eta$) eines Tripels in ein Tripel über, und daraus lässt sich beweisen,
dass jede Gleichung 9$^{\mathrm{ten}}$ Grades, deren Galois'sche Gruppe $P$ in der
vollständigen linearen Congruenzgruppe enthalten ist, eine Tripelgleichung ist, wenn
noch die weitere Bedingung hinzukommt, dass durch die Gruppe $P$ irgend zwei gegebene
Wurzeln in zwei andere gleichfalls beliebig gegebene Wurzeln übergehen, oder, was damit
gleichbedeutend ist, wenn $P$ \gesperrt{zweifach transitiv} ist.

Denn wenn $P$ aus lauter linearen Substitutionen besteht, so reducirt sich $P$ durch
Adjunction zweier beliebiger Wurzeln $x_1 = (\xi_1, \eta_1)$, $x_2 = (\xi_2, \eta_2)$
auf eine Gruppe $P_1$, die nicht nur $x_1, x_2$, sondern auch die dritte Wurzel
$x_3 = (\xi_3, \eta_3)$ des durch $x_1, x_2$ bestimmten Tripels ungeändert lässt, und
folglich gehört $x_3$ dem neuen Rationalitätsbereiche an, und mithin ist $x_3$ rational
durch $x_1, x_2$ ausdrückbar in der Form

\weq{8}{x_3 = \Theta (x_1, x_2).}

Wenn nun die Gruppe $P$ zweifach transitiv ist, so giebt es eine Permutation in $P$,
durch die $x_1, x_2$ in zwei beliebige andere Wurzeln $y_1, y_2$ übergehen, und durch
die folglich auch $x_3$ in die dritte Wurzel des Tripels $y_1, y_2, y_3$ übergeht, und
diese Permutation ist auf (8) anwendbar. Also ist auch

\[
y_3 = \Theta (y_1, y_2).
\]

Wir sprechen dies als Satz aus:

3. \gesperrt{Jede Gleichung 9$^{\mathrm{ten}}$ Grades, deren Gruppe linear
und zweifach transitiv ist, ist eine Tripelgleichung.}

Durch die Permutationen einer zweifach transitiven linearen Gruppe kann jedes Tripel
in jedes andere, und zwar in beliebiger Anordnung, übergeführt werden.

Die Voraussetzung der zweifachen Transitivität kann auch so ausgedrückt werden, dass
alle die Permutationen aus $P$, die eine Wurzel ungeändert lassen, die übrigen noch
transitiv mit einander verbinden. Bezeichnen wir mit $P_0$ die in $P$ enthaltene
Gruppe, durch deren Permutationen eine der Wurzeln $x_0$ ungeändert bleibt, so muss
$P_0$ in Bezug auf die anderen Wurzeln noch transitiv sein. Ist aber die Gruppe $P$
nur einfach transitiv,
%CONT>>
